% short/00_abstract.tex: the abstract (a plain summary; the only part allowed to use terms explained later).
% Numbers only through \R keys (tables/numbers_index.tsv). The three-attempt accuracy is final/att_k3_used_acc, the
% combination the model uses and the value plotted in the attempts figure; it is quoted with its own one-attempt
% value (final/att_k1_used_acc, the same test) so that it does not read as a rise from the honest estimate.
\begin{abstract}
\noindent A person brings a smartphone from arm's length to the nose with the left or the right hand. We asked
whether the phone's motion sensors alone can tell which hand it was, as a probability \pright{} that the right
hand moved the phone, for people the model has never seen. We built a web app that records the movement and
collected \R{describe/primary_trials} usable movements from \R{describe/n_persons} participants. We tested several
kinds of models, always on a participant left out of training, and a one-dimensional convolutional neural network
(1D CNN) performed best: for new participants we estimate that it identifies the hand with an accuracy of
\R{final/nested_accuracy} (95\,\% confidence interval \Rb{final/nested_accuracy_lo}--\R{final/nested_accuracy_hi}),
with an area under the ROC curve of \R{final/nested_auc}. We also checked that the predicted probabilities can be taken at face value (calibration):
some models were over-confident and needed a correcting function, while the network's outputs were already close to
calibrated. Finally, combining several movements made with the same hand improves the prediction: in a test of this
step, the accuracy rose from \R{final/att_k1_used_acc} with one movement to \R{final/att_k3_used_acc} with three.
\end{abstract}
